\selectlanguage{ngerman}

\maketitle

\thispagestyle{empty}
\vspace*{2.25cm}
$$\frac{1}{\pi} = 12\cdot\sum_{n=0}^\infty \frac{(-1)^n(6n)!}{(3n)!(n!)^3}\cdot \frac{13591409 + 545140134 n}{ 640320^{3n+3/2}}$$
\vspace*{2.25cm}
\label{\en{EN}titlepage}

{\narrower\narrower\narrower{
\noindent{\scshape Abstract.} In this paper we give another proof of the Chudnovsky formula for calculating $\pi$ -- a proof in detail with means of basic complex analysis.

With the exception of the tenth chapter, the proof is self-contained, with proofs provided for all the advanced theorems we use (e.g.~for the Clausen formula and for the Picard-Fuchs differential equation).


\vspace*{2mm}
{\begin{center}\emph{English version: pp.~\pageref{ENtitlepage}--\pageref{ENliterat}}\end{center}}

\vspace*{1.5cm}

\noindent{\scshape Zusammenfassung.} In diesem Aufsatz wird die Chudnovsky-Formel zur Berechnung von $\pi$ erneut bewiesen -- wesentlich ausführlicher, mit elementaren Methoden der Funktionentheorie und der Analysis.

Die benötigten fortgeschrittenen Sätze (z.B.~die Clausen-Formel und die Picard-Fuchs-Differentialgleichung) werden ihrerseits ausführlich bewiesen. Nur im zehnten Kapitel verweisen wir auf externe Quellen.


\vspace*{2mm}
{\begin{center}\emph{Deutsche Version: S.~\pageref{titlepage}--\pageref{literat}}\end{center}}

~
}}
